\subsection{O relógio de uma ligação}
\label{subsec:dandelion-fila-ligacao}
\index{dandelion++@Dandelion++!filas de difusao@filas de difusão}

A transacção de Alice alcançou o nó difusor. O que significa agora a
expressão «enviar aos pares» para os vizinhos que aguardam o pagamento?
Consideremos duas ligações públicas elegíveis, sempre com os nomes
referidos ao difusor. A ligação a Bruno foi iniciada pelo difusor e é,
portanto, de saída. A ligação de Clara foi iniciada por Clara e é de
entrada. Os dados da transacção vão do difusor para o vizinho em ambos
os casos. É a iniciativa de estabelecer a ligação que dá nome ao
relógio escolhido.

Partamos de duas filas vazias no instante zero. A cada uma associa-se
um prazo quando chega o primeiro elemento. Suponhamos que saíram dez
quartos de segundo para Bruno e vinte e quatro para Clara. São
resultados possíveis das distribuições atrás indicadas. As filas
abrem com a transacção \(T\) de Alice e têm os envios previstos para
\(2{,}5\) e \(6\) segundos, respectivamente. No instante
\(1{,}5\) chega outra transacção \(U\), que o difusor também admitiu
para ambos os destinos. As duas filas já possuem um prazo. Cada uma
recebe mais um elemento e conserva o instante de envio que tinha.

Para acompanhar o exemplo, admitamos que cada envio se realiza no
instante previsto e que a entrega é imediata. O tempo zero é o da
chegada de \(T\) ao difusor. Uma terceira transacção, \(W\), fica
disponível no instante 4. O quadro acompanha apenas estas ligações e
admite que os três corpos vieram de outros pares.
\begin{center}
\begin{tabular}{r|l|l}
instante (s)&fila para Bruno&fila para Clara\\
\hline
\(0\)&\(T\), envio a \(2{,}5\)&\(T\), envio a 6\\
\(1{,}5\)&\(T,U\), envio a \(2{,}5\)&\(T,U\), envio a 6\\
\(2{,}5\)&saem \(T,U\); fila vazia&\(T,U\), envio a 6\\
\(4\)&\(W\), novo envio a 7&\(T,U,W\), envio a 6\\
\(6\)&\(W\), envio a 7&saem \(T,U,W\); fila vazia\\
\(7\)&sai \(W\); fila vazia&vazia
\end{tabular}
\end{center}
Para a segunda abertura da fila de Bruno escolhemos doze quartos de
segundo. O prazo ficou em \(4+3=7\). A transacção \(W\), que chegou
depois das outras duas, alcançou primeiro Clara; \(T\) e \(U\) tinham
alcançado primeiro Bruno. A lista de letras do quadro representa os
corpos presentes na fila. Na formação do lote os corpos são ordenados
e eliminam-se os duplicados \cite{monero-source-v16}.

Na fila para Clara, \(T\) esperou seis segundos, \(U\) esperou
\(6-1{,}5=4{,}5\) segundos e \(W\) esperou apenas dois. Houve
uma escolha aleatória de prazo para todo esse lote. Na fila para Bruno
houve duas escolhas, pois a fila esteve vazia entre \(2{,}5\) e 4.
Assim se lê correctamente a média de cinco segundos da distribuição
de entrada. Ela vale para um novo prazo, contado a partir da abertura
da fila; o corpo que entra numa fila já aberta aguarda o tempo que
resta. A distinção faz diferença mesmo quando os relógios funcionam
perfeitamente e o tempo de transmissão é desprezável.

\paragraph{Que unidade tem a variância?}
As médias das duas distribuições permitem também uma pequena conta
sobre a dispersão dos prazos. Seja \(X\) a duração em quartos de
segundo, com distribuição de Poisson de média \(\lambda\). Nessa
unidade \(\operatorname{Var}(X)=\lambda\). Para obter segundos,
multiplicamos cada duração por \(1/4\). A média fica multiplicada
por \(1/4\) e a variância por \(1/16\), pois a diferença à média é
elevada ao quadrado. Para a ligação de entrada resulta
\[
  \mathbb E[D_{\rm in}]=20/4=5\ {\rm s},\qquad
  \operatorname{Var}(D_{\rm in})=20/16=1{,}25\ {\rm s}^2 .
\]
O desvio padrão é \(\sqrt{1{,}25}\simeq1{,}12\) segundos. Na de
saída a mesma conta dá \(10/4=2{,}5\) segundos de média e
\(\sqrt{10/16}\simeq0{,}79\) segundos de desvio padrão. Uma
distribuição de Poisson gerada directamente em segundos com média 5
teria variância \(5\ {\rm s}^2\). Escolher a unidade da amostra
altera, pois, a dispersão dos atrasos físicos mesmo quando se conserva
o mesmo tempo médio.

Voltemos à carteira de Bob, que já conhece os dados da saída que espera
receber. Se o seu nó for um dos pares que recebe o corpo, pode começar
a examiná-lo logo depois do envio e da recepção. Outro nó pode receber
esse mesmo corpo num lote posterior, ainda que se encontre à mesma
distância do difusor. Ambos têm então a possibilidade de validar os
mesmos dados da transacção. A altura a que um bloco vier a incluí-la
é uma informação posterior, dada pelo ramo aceite. O relógio das filas
explica a diferença entre notícias que chegam a dois vizinhos; o
número de confirmações começa a contar numa escala que os cabeçalhos
dos blocos tornam comum.
